\documentclass[11pt]{article}
\usepackage[margin=1in]{geometry}
\usepackage{amsmath,amssymb,amsthm}
\usepackage{booktabs}
\usepackage{array}
\usepackage{microtype}
\usepackage{enumitem}
\usepackage[colorlinks=true,linkcolor=blue,citecolor=blue,urlcolor=blue]{hyperref}
\usepackage{xcolor}

\newtheorem{lemma}{Lemma}
\newtheorem{proposition}{Proposition}
\newtheorem{remark}{Remark}
\newcommand{\R}{\mathbb{R}}
\newcommand{\E}{\mathbb{E}}
\newcommand{\norm}[1]{\left\lVert #1 \right\rVert}
\newcommand{\sg}{\mathrm{sg}}

\title{\textbf{FE-JEPA: Energy-Anchored Joint-Embedding Predictive Pretraining
for Finite-Element Surrogates}\\[1ex]
\large A Feasibility Study, State-of-the-Art Survey, and Research Execution Plan}
\author{Prepared for Ruifeng Cao}
\date{June 12, 2026}

\begin{document}
\maketitle

\begin{abstract}
\noindent
This document examines the feasibility of combining Joint-Embedding Predictive
Architectures (JEPA) with physics-informed learning for finite-element analysis
(FEA), surveys the state of the art as of June 2026, and proposes a concrete
research topic with a detailed execution plan. The intersection has just become
an active frontier: \emph{PI-JEPA} (April~2026) applies physics-regularized
masked latent prediction to subsurface flow on regular grids, and
\emph{AeroJEPA} (May~2026) learns JEPA latents of 3D aerodynamic fields.
No published work, to the best of this survey's knowledge, applies JEPA to
unstructured-mesh structural FEA, and the only physics-anchored JEPA reports
that its strong-form finite-difference residual was \emph{neutral or mildly
harmful}. We argue that solid mechanics is precisely the setting where physics
anchoring should work, because the discrete \emph{variational} energy
$\Pi_h(\hat u)=\tfrac12\hat u^\top K\hat u - F^\top\hat u$, computable by cheap
FE assembly without any solve, is---exactly, not approximately---the
energy-norm distance to the FE solution up to a constant
(Lemma~\ref{lem:energy}). The proposed topic, \textbf{FE-JEPA}, is a label-free
pretraining framework for FEA surrogates on unstructured meshes that combines
(i) masked geometry/boundary-condition latent prediction, (ii) mesh-refinement
views as \emph{exact} physical augmentations, (iii) an assembled-energy anchor
whose gradient coincides with the supervised energy-norm gradient, and (iv)
provable collapse control via isotropic-Gaussian regularization (LeJEPA/SIGReg).
We give the formal setup, two theory targets with one complete proof, five
research questions, falsification-first experiments, baselines, datasets,
metrics, a 12-month phased plan with go/no-go gates, a compute budget, and a
risk register. Overall feasibility is assessed as \emph{high} on engineering
grounds and \emph{medium-high} on novelty grounds, with speed of execution the
dominant strategic risk.
\end{abstract}

\tableofcontents
\bigskip

\section{Introduction and Motivation}
\label{sec:intro}

Three independently mature research programs meet in this proposal.

\paragraph{FEA surrogates.}
Finite-element analysis is the workhorse of structural engineering, but
high-fidelity solves are too expensive for the many-query regimes that matter
in practice: design-space exploration, topology and shape optimization,
uncertainty quantification, real-time monitoring, and digital twins. A decade
of work on learned surrogates---convolutional models on rasterized geometry,
graph networks on meshes \cite{pfaff2021}, neural operators
\cite{li2021fno,lu2021deeponet,kovachki2023,li2023geofno}, and
geometry-general transformers \cite{wu2024transolver}---has produced fast
approximations of displacement and stress fields. The persistent bottleneck is
\emph{labeled data}: every training pair requires a full FE solve, and
accuracy degrades sharply out of the training distribution of geometries,
loads, and boundary conditions. Benchmark efforts such as SimJEB
\cite{whalen2021simjeb} were created precisely because realistic geometric
diversity overwhelms surrogates trained on synthetic families.

\paragraph{Physics-informed learning.}
PINNs \cite{raissi2019} inject the governing equations into the loss and in
principle remove the need for labels. In computational solid mechanics they
have spawned a large literature---strong-form collocation, variational/energy
forms (Deep Ritz \cite{e2018ritz}, the Deep Energy Method
\cite{samaniego2020,nguyen2020dem}), and FEM--network hybrids such as I-FENN
\cite{pantidis2023ifenn} and deep-finite-element couplings---surveyed at
length in a 2026 review of PINNs for solid and structural
mechanics.\footnote{\emph{Physics-Informed Neural Networks: Current Progress
and Challenges in Computational Solid and Structural Mechanics}, CMES, vol.~146,
no.~2, 2026. \url{https://www.techscience.com/CMES/v146n2/66323/html}}
Yet PINNs remain predominantly \emph{per-instance} solvers: each new geometry
or load requires fresh optimization, which is typically slower and less
robust than a direct FE solve for the same problem; well-documented failure
modes of strong-form residual training \cite{krishnapriyan2021,wang2022ntk}
compound the issue. Physics-informed \emph{operator} learning (PINO
\cite{li2024pino}, physics-informed DeepONets) amortizes across instances but
still anchors physics through continuous-form residuals that are only
approximately consistent with the discrete solver one ultimately wants to
match.

\paragraph{JEPA.}
Joint-Embedding Predictive Architectures \cite{lecun2022} learn by predicting
the \emph{latent representation} of masked or future content rather than
reconstructing raw signals, and have progressed from images (I-JEPA
\cite{assran2023ijepa}) to video world models (V-JEPA \cite{bardes2024vjepa},
V-JEPA~2 \cite{meta2025vjepa2}). Two recent theoretical advances make JEPA
attractive as a \emph{principled} component rather than a bag of heuristics:
LeJEPA \cite{balestriero2025lejepa} proves that isotropic Gaussian embeddings
minimize worst-case downstream risk and enforces them with a linear-time
sketched regularizer (SIGReg) that replaces stop-gradients, EMA teachers, and
schedule tricks; and follow-up work gives identifiability conditions under
which a LeJEPA-style learner provably recovers latent world state
\cite{klindt2026}.

\paragraph{Why combine them, and why for FEA specifically.}
The economics of FEA exhibit exactly the \emph{data asymmetry} that motivates
JEPA-style pretraining for simulation, as articulated by PI-JEPA
\cite{yee2026pijepa}: unlabeled \emph{inputs} (CAD geometries, meshes,
material fields, load cases) are abundant and nearly free, while labeled
\emph{outputs} (solved fields) are expensive. JEPA is the natural mechanism
for exploiting unlabeled inputs, since it never needs the output field to
define its pretraining signal. What JEPA alone cannot do is connect latents to
the \emph{solution}; that is what a physics term must supply. The central
technical observation of this proposal is that linear(ized) solid mechanics
offers a physics term of unusual quality: the assembled discrete energy
functional is cheap to evaluate (sparse matrix--vector products, no solve), is
\emph{exactly} minimized by the FE solution, and---crucially---its gradient
with respect to a candidate field equals the gradient of the supervised
energy-norm loss against the unknown FE solution (Section~\ref{sec:theory}).
Label-free training and supervised training literally coincide at the level of
gradients. No strong-form, autodiff-based PINN residual has this property, and
the one existing physics-anchored JEPA, which used finite-difference
strong-form residuals on grids, reported the physics term to be neutral or
mildly harmful \cite{yee2026pijepa}. FEA is therefore not merely an
application domain for ``JEPA + PINN''; it is the domain where the combination
can be made \emph{exact}, and where the open question left by PI-JEPA---does a
\emph{discretization-consistent} physics anchor help where a
discretization-inconsistent one did not?---can be answered cleanly.

\section{State of the Art (as of June 2026)}
\label{sec:sota}

\subsection{Learned surrogates for FEA}

Mesh-based graph networks remain the default for unstructured FEA.
MeshGraphNets \cite{pfaff2021} established message passing on simulation
meshes (including the quasi-static \textsc{DeformingPlate} hyperelastic
benchmark), and a stream of derivatives applies the recipe to stress-field
prediction on plates with holes, brackets, femurs, and multi-component
assemblies; representative recent work includes mesh-based geometric deep
learning for large multi-component structures (CMAME, 2026). Neural operators
provide the resolution-robust alternative: FNO \cite{li2021fno} on regular
grids, Geo-FNO \cite{li2023geofno} for general geometries via learned
deformations (whose public elasticity benchmark is a standard testbed),
DeepONet \cite{lu2021deeponet}, and the general operator framework of
\cite{kovachki2023}. Transformer solvers such as Transolver
\cite{wu2024transolver} currently set strong marks on geometry-general PDE
benchmarks by attending over learned physical slices rather than raw mesh
points. On the benchmark side, SimJEB \cite{whalen2021simjeb} provides 381
human-designed jet-engine brackets with displacement and von~Mises stress
fields under four load cases, explicitly framed as a surrogate-modeling
challenge because its geometric diversity exceeds synthetic datasets; PDEBench
\cite{takamoto2022pdebench} covers grid-based PDE suites. The shared weakness
across this entire line: training is fully supervised on solved fields, and
label efficiency is rarely the headline metric.

\subsection{Physics-informed learning in solid mechanics}

Strong-form PINNs \cite{raissi2019} are well studied for elasticity,
elastodynamics, and fracture, but suffer documented optimization pathologies
\cite{krishnapriyan2021,wang2022ntk} and remain per-instance. The variational
family is the more natural fit for solids: Deep Ritz \cite{e2018ritz}
minimizes the energy functional directly; the Deep Energy Method
\cite{samaniego2020,nguyen2020dem} does so for finite-deformation
hyperelasticity, where a stored-energy potential exists. Hybrid lines couple
networks with FEM rather than replacing it: I-FENN embeds networks inside FE
assembly for nonlocal damage \cite{pantidis2023ifenn}; ``deep finite element''
frameworks and transfer-learning-accelerated FEM--PINN couplings (CMAME 2025;
IJMS 2025) report accurate static response for complex shells while
acknowledging that low efficiency and precision still hinder plain PINNs on
boundary-value problems. The 2026 CMES review (footnote, p.~2) confirms the
field's own diagnosis: physics-informed losses are valuable as
\emph{regularization and label replacement}, but per-instance physics-only
optimization does not displace FEM. Physics-informed operator learning (PINO
\cite{li2024pino}) is the amortized compromise, yet PI-JEPA's controlled
comparison found PINO roughly on par with plain FNO in the low-label regime it
studied \cite{yee2026pijepa}---evidence that \emph{how} physics is injected
matters more than \emph{whether}.

\subsection{JEPA: methods and theory}

I-JEPA \cite{assran2023ijepa} predicts masked image-patch embeddings; V-JEPA
\cite{bardes2024vjepa} and V-JEPA~2 \cite{meta2025vjepa2} extend to video and
action-conditioned planning. The framework has diffused rapidly into
specialized domains (EEG, ultrasound, genomics, point clouds, time series,
graphs), indicating that the recipe transfers when the masking structure is
adapted to the modality. Two results change the design space for scientific
applications. First, LeJEPA \cite{balestriero2025lejepa} proves the isotropic
Gaussian is the embedding distribution minimizing worst-case downstream
prediction risk and enforces it via SIGReg---random 1D projections matched to
a Gaussian by a characteristic-function test---yielding a single-hyperparameter,
linear-time objective (about 50 lines of code) that removes EMA teachers and
stop-gradient asymmetries while reaching 79\% ImageNet top-1 with ViT-H/14.
Second, identifiability analysis \cite{klindt2026} shows LeJEPA-style training
recovers latent world state (up to linear maps) if and only if latents are
isotropic Gaussian with stationary additive-noise transitions---and,
instructively, that goal-directed (non-isotropic) data sampling breaks
recovery. For a proposal whose scientific claim is ``the latent space encodes
the mechanics,'' these are the right foundations: collapse control becomes a
testable statistical condition rather than a trick, and the data-sampling
caveat directly informs how pretraining distributions over loads and
geometries should be designed.

\subsection{The intersection: latent prediction meets physics (2025--2026)}
\label{sec:intersection}

The specific combination of JEPA-style latent prediction with
physics-informed objectives crystallized only in the last six months.

\paragraph{PI-JEPA (arXiv:2604.01349, April 2026) \cite{yee2026pijepa}.}
The closest existing work. Setting: subsurface flow (two-phase Darcy,
reactive transport) on \emph{regular grids}. Method: masked latent prediction
on \emph{unlabeled input parameter fields} (permeability, porosity) with a
Fourier-enhanced patch encoder, an EMA target encoder plus an explicit
collapse-prevention regularizer, a bank of latent predictors structurally
aligned with the Lie--Trotter operator splitting of the governing equations,
and per-sub-operator \emph{strong-form finite-difference} PDE residuals as
physics regularization; fine-tuning with as few as 100 labeled runs. Reported
findings are unusually candid and shape this proposal: (i) self-supervised
pretraining contributes a real but moderate 6--14\% improvement; (ii) the
operator-split \emph{architecture}, not pretraining, is the dominant factor;
(iii) the PDE residual term is \emph{neutral}---ablating it slightly improved
accuracy ($-8.1\%$ error)---which the authors attribute to the
finite-difference approximation; (iv) extension to non-elliptic systems and
other splittings is posed as open. Items (iii)--(iv) define the gap this
proposal attacks: replace an approximation-quality residual with a
discretization-exact variational one, in a domain (elliptic solid statics)
the PI-JEPA benchmarks did not touch.

\paragraph{AeroJEPA (arXiv:2605.05586, May 2026).}
\footnote{Author list not recorded in this survey; verify on arXiv before
citing. \url{https://arxiv.org/abs/2605.05586}}JEPA latents for 3D aerodynamic \emph{solution fields}: context/target
encoders over CFD fields, a latent predictor, and an implicit-neural-
representation decoder, evaluated on high-lift and transonic-wing datasets
with latent-space design optimization. Differences from this proposal: it
pretrains on \emph{labeled} solution fields (the costly resource), targets
external aerodynamics rather than structural FEA, and uses no physics term.

\paragraph{Latent dynamics and generative solvers.}
Latent Generative Solvers (arXiv:2602.11229, Feb.~2026)\footnote{Author list
not recorded in this survey. \url{https://arxiv.org/abs/2602.11229}} pair a
pretrained VAE latent space with flow-matched transformer dynamics for
long-horizon multi-PDE rollout---reconstruction-based rather than JEPA, and
time-dependent rather than static, but evidence that shared latent physics
spaces are the current frontier.

\paragraph{Geometry/structure self-supervision without physics.}
The ``Shape'' industrial-CAD foundation model (arXiv:2604.22826,
April~2026)\footnote{Author list not recorded in this survey.
\url{https://arxiv.org/abs/2604.22826}} pretrains a 10.9M-parameter backbone
on 61k CAD meshes (Thingi10K, MFCAD, Fusion360) via masked-token
reconstruction of geometry statistics---demonstrating that exactly the
unlabeled corpus FE-JEPA needs already exists at scale, but using
reconstruction rather than latent prediction and no mechanics. A polycrystal
microstructure foundation model (arXiv:2512.06770, Dec.~2025) similarly shows
masked 3D pretraining improving homogenized-stiffness prediction.

\paragraph{PDE foundation models (pre-JEPA lineage).}
Multiple Physics Pretraining \cite{mccabe2023mpp}, Poseidon
\cite{herde2024poseidon}, and DPOT \cite{hao2024dpot} pretrain on large
corpora of \emph{solved trajectories} with supervised or autoregressive
denoising objectives. They establish that pretraining transfers across PDE
families, but their pretraining currency is exactly the expensive resource
(solutions) that JEPA-style input-side pretraining avoids spending.

\section{Gap Analysis and Feasibility Assessment}
\label{sec:gap}

\subsection{Gap matrix}

Table~\ref{tab:gap} positions the proposal against the closest work along the
axes that define it.

\begin{table}[h]
\centering
\footnotesize
\setlength{\tabcolsep}{3.5pt}
\begin{tabular}{@{}p{3.0cm}cccccc@{}}
\toprule
 & \textbf{Domain} & \textbf{Unstr.} & \textbf{Label-free} & \textbf{Physics} & \textbf{Physics form} & \textbf{Collapse}\\
 & & \textbf{mesh} & \textbf{pretrain} & \textbf{anchor} & & \textbf{control}\\
\midrule
PI-JEPA \cite{yee2026pijepa} & subsurface flow & no (grid) & yes & yes & strong-form FD & EMA + reg.\\
AeroJEPA (2026) & aerodynamics & point cloud & no (fields) & no & --- & EMA\\
Shape (2026) & CAD geometry & yes & yes & no & --- & (recon.)\\
MPP / Poseidon / DPOT & multi-PDE & no & no & no & --- & (recon./AR)\\
PINO \cite{li2024pino} & multi-PDE & no & partial & yes & strong-form & ---\\
Deep Ritz / DEM \cite{e2018ritz,samaniego2020} & solids & n/a & yes & yes & variational & --- (per inst.)\\
\textbf{FE-JEPA (proposed)} & \textbf{structural FEA} & \textbf{yes} & \textbf{yes} & \textbf{yes} & \textbf{assembled energy} & \textbf{SIGReg}\\
\bottomrule
\end{tabular}
\caption{Gap matrix. ``Label-free pretrain'' = pretraining signal requires no
solved fields. ``Assembled energy'' = discrete variational functional built by
FE assembly on the instance's own mesh, exact w.r.t.\ the FE solution.}
\label{tab:gap}
\end{table}

No surveyed work occupies the bottom row. Targeted searches (June 2026) for
JEPA applied to finite elements, solid mechanics, or structural analysis
return PI-JEPA and AeroJEPA as the only physics-simulation JEPAs, neither in
structural FEA, neither on unstructured FE meshes, and neither with a
discretization-exact physics term.

\subsection{Feasibility verdict}

\paragraph{Engineering feasibility: high.}
Every ingredient is off the shelf. Unlabeled geometry corpora exist at the
required scale (Thingi10K/MFCAD/Fusion360 as used by Shape; ABC-style CAD
collections; parametric generators). Meshing and assembly are commodity
operations (Gmsh; FEniCSx, scikit-fem, or CalculiX for $K,F$ assembly);
assembly without solving is $O(\text{nnz})$ and embarrassingly parallel, so
the energy anchor adds only sparse mat-vecs per training step. Labeled
evaluation sets exist (SimJEB; Geo-FNO elasticity; MeshGraphNets
DeformingPlate; in-house ABAQUS/CalculiX generation for controlled studies).
Encoders (graph transformer / Transolver-class), JEPA training loops, and
SIGReg (public reference implementation) are all published. Nothing requires
new infrastructure beyond careful integration.

\paragraph{Scientific feasibility: high, with one honest caveat.}
The core lemma (Section~\ref{sec:theory}) guarantees the physics anchor is
\emph{informative}---it is the supervised energy-norm loss in disguise---so
the failure mode that PI-JEPA observed for finite-difference residuals is
structurally excluded for linear problems. The caveat: PI-JEPA also found
that \emph{architecture} dominated \emph{pretraining} in its setting. It is
entirely possible that the same holds here and the JEPA component contributes
single-digit percentages while the energy anchor does the heavy lifting. That
outcome is still a publishable, decision-relevant finding (it would say
``amortized Ritz training, not SSL, is what FEA surrogates were missing''),
but the proposal's experimental design must be able to attribute credit
cleanly (Section~\ref{sec:plan}).

\paragraph{Novelty feasibility: medium-high, time-sensitive.}
Two adjacent papers appeared within the last 90 days; PI-JEPA's own
future-work section gestures toward other operator classes; AeroJEPA's
machinery (INR decoder over JEPA latents) is one dataset away from solid
mechanics. The window for ``first JEPA for structural FEA with an exact
variational anchor'' is open now and should be assumed to close within
6--12 months. The differentiators that survive even if scooped on the
application: the exactness lemma and gradient identity, the
mesh-refinement-as-exact-augmentation principle, and the
identifiability-style anti-collapse analysis---these are \emph{theory}
contributions a rushed applications paper is unlikely to include.

\section{Proposed Research Topic: FE-JEPA}
\label{sec:proposal}

\subsection{One-sentence statement}

\emph{Label-free pretraining of FEA surrogates on unstructured meshes by
joint-embedding prediction over masked geometry and boundary conditions,
anchored to the mechanics through the assembled discrete energy functional,
with provable collapse control and provable equivalence---for linear
problems---between the label-free anchor and supervised energy-norm
training.}

\subsection{Formal setup}
\label{sec:setup}

Consider a family of linear elastostatic boundary-value problems indexed by
$\theta=(\Omega,\,E,\nu,\,\Gamma_D,\Gamma_N,\,f,\,t)$: domain, material
parameters, Dirichlet/Neumann partition, body force, traction. Weak form:
find $u\in V$ with $a_\theta(u,v)=\ell_\theta(v)$ for all $v\in V$, where
$a_\theta(u,v)=\int_\Omega \sigma(u):\varepsilon(v)\,dx$ and
$\ell_\theta(v)=\int_\Omega f\cdot v\,dx+\int_{\Gamma_N} t\cdot v\,ds$.
A mesh $\mathcal{T}_h(\theta)$ and FE space $V_h$ give the discrete system
$K\,U^\star=F$ with $K$ symmetric positive definite after Dirichlet
elimination, and the discrete total potential energy
\begin{equation}
\Pi_h(\hat u)\;=\;\tfrac12\,\hat u^\top K\,\hat u\;-\;F^\top \hat u,
\qquad \hat u\in\R^{n}.
\label{eq:energy}
\end{equation}
Assembling $(K,F)$ costs one pass over elements; \emph{no linear solve is
performed at training time}. The unlabeled pretraining corpus is a
distribution $\mathcal{D}$ over instances $\theta$ (geometries $\times$
load/BC programs $\times$ material draws), each meshed at one or more
resolutions.

\paragraph{Architecture.}
A shared encoder $f_\xi$ maps tokenized instance data---node coordinates,
element connectivity, material fields, BC/load tokens---to per-region latents
$z$ (graph-transformer or Transolver-class backbone, so that token count
decouples from node count). A predictor $p_\phi(z_{\mathrm{ctx}},m)$ predicts
target-region latents from context latents and a mask/condition descriptor
$m$ (which carries the JEPA latent variable: load case, masked-region
location, mesh-resolution tag). A lightweight field decoder $g_\psi$ maps
latents to a nodal displacement field $\hat u=g_\psi(z)\in\R^{n}$ on the
instance's own mesh (per-node head conditioned on pooled/region latents, so
$g_\psi$ is mesh-size-agnostic).

\subsection{Training objectives}
\label{sec:objectives}

\begin{align}
\mathcal{L}
&= \underbrace{\E_{\theta,\,\text{mask}}\,
   \norm{\,p_\phi(z_{\mathrm{ctx}},m)-z_{\mathrm{tgt}}\,}_2^2}_{\mathcal{L}_{\mathrm{pred}}\;:\;\text{masked latent prediction}}
\;+\;\lambda_E\,
\underbrace{\E_{\theta,\,j}\;\Pi_h^{(\theta,j)}\!\bigl(g_\psi(z)\bigr)}_{\mathcal{L}_{\mathrm{phys}}\;:\;\text{assembled-energy anchor}}
\notag\\[0.5ex]
&\qquad +\;\lambda_S\,
\underbrace{\mathrm{SIGReg}(z)}_{\text{collapse control}}
\;+\;\lambda_I\,
\underbrace{\E_{\theta}\,
   \norm{\,h(z^{(h_1)})-h(z^{(h_2)})\,}_2^2}_{\mathcal{L}_{\mathrm{inv}}\;:\;\text{cross-mesh invariance}}.
\label{eq:loss}
\end{align}

Here $j$ ranges over a small battery of load cases applied to the same
geometry (cheap: $F$ changes, $K$ is reused), $z^{(h_1)},z^{(h_2)}$ are
latents of the \emph{same} instance meshed at two resolutions, and $h$ is a
projection head. Three design choices deserve emphasis.

\paragraph{(a) The anchor is exact, not a PINN proxy.}
$\mathcal{L}_{\mathrm{phys}}$ uses the instance's own assembled $(K,F)$.
Its minimizer over $\hat u$ is the FE solution itself, so the surrogate's
fixed point is the object FEA users actually want---not a continuous-PDE
stand-in whose discretization mismatch PI-JEPA identified as the likely
culprit for its neutral physics term. Section~\ref{sec:theory} makes this
exactness quantitative.

\paragraph{(b) Mesh refinement is an \emph{exact} augmentation.}
In vision, augmentations (crops, colors) only approximately preserve
semantics. Here, two meshes of the same boundary-value problem represent
\emph{identical} physics by construction; $\mathcal{L}_{\mathrm{inv}}$
therefore demands discretization-invariant latents with zero augmentation
bias---a property neural-operator practice prizes
\cite{kovachki2023,li2023geofno} but SSL has not exploited. This is, to this
survey's knowledge, a new augmentation principle for scientific SSL.

\paragraph{(c) Collapse control is statistical, not heuristic.}
SIGReg \cite{balestriero2025lejepa} replaces EMA/stop-gradient machinery,
keeps the objective symmetric (important because context and target here are
the same modality), and---via \cite{klindt2026}---turns ``the latents encode
the mechanics'' into a testable distributional condition. The energy anchor
itself adds a second, physics-grounded anti-collapse force, formalized as
Proposition~\ref{prop:collapse}.

\subsection{Theory targets}
\label{sec:theory}

\begin{lemma}[Exactness of the assembled-energy anchor]
\label{lem:energy}
Let $K\in\R^{n\times n}$ be symmetric positive definite, $F\in\R^n$,
$U^\star=K^{-1}F$, and $\Pi_h$ as in \eqref{eq:energy}. Then for every
$\hat u\in\R^n$,
\[
\Pi_h(\hat u)-\Pi_h(U^\star)\;=\;\tfrac12\,\norm{\hat u-U^\star}_K^2,
\qquad\text{and}\qquad
\nabla_{\hat u}\,\Pi_h(\hat u)\;=\;K(\hat u-U^\star)
\;=\;\nabla_{\hat u}\,\tfrac12\norm{\hat u-U^\star}_K^2 .
\]
\end{lemma}

\begin{proof}
$\Pi_h(\hat u)-\Pi_h(U^\star)
=\tfrac12\hat u^\top K\hat u-F^\top\hat u
+\tfrac12 U^{\star\top}KU^\star$ using
$\Pi_h(U^\star)=-\tfrac12 U^{\star\top}KU^\star$ (since $KU^\star=F$).
Substituting $F=KU^\star$ in $-F^\top\hat u$ and completing the square gives
$\tfrac12(\hat u-U^\star)^\top K(\hat u-U^\star)$. The gradient identity is
immediate: $\nabla\Pi_h=K\hat u-F=K(\hat u-U^\star)$.
\end{proof}

\begin{remark}
The consequence is stronger than ``same minimizer'': the \emph{label-free}
loss $\mathcal{L}_{\mathrm{phys}}$ and the \emph{supervised} energy-norm loss
produce identical gradients with respect to the decoded field, for every
instance, at every training step, without ever computing $U^\star$. By
Galerkin coercivity the energy norm controls the $H^1$ error, so this is
supervision in the norm FE theory cares about. Per-instance constants
$\Pi_h(U^\star)$ are unknown during training (irrelevant to gradients) and
are computed only on the labeled evaluation split to report the
\emph{energy gap} metric. For finite-deformation hyperelasticity, $\Pi_h$
remains a well-defined (nonconvex) stored-energy functional and the anchor
generalizes with stationarity replacing global minimality; for plasticity and
contact, incremental variational formulations provide the analogue and are
deferred to a stretch phase.
\end{remark}

\begin{proposition}[Energy anchor prevents collapse; to be formalized]
\label{prop:collapse}
Fix a geometry and loads $F_1,\dots,F_m$ spanning an $r$-dimensional
subspace, with solutions $U_j^\star=K^{-1}F_j$ spanning an $r$-dimensional
subspace of $\R^n$. If the decoder is $L$-Lipschitz (latent metric to
$K$-norm) and each instance's decoded field achieves energy gap
$\Pi_h(g_\psi(z_j))-\Pi_h(U_j^\star)\le\varepsilon$, then by
Lemma~\ref{lem:energy} $\norm{g_\psi(z_j)-U^\star_j}_K\le\sqrt{2\varepsilon}$,
hence pairwise latent distances obey
$\norm{z_i-z_j}\ \ge\ \bigl(\norm{U_i^\star-U_j^\star}_K-2\sqrt{2\varepsilon}\bigr)/L$:
latents of mechanically distinct instances cannot collapse, with separation
quantified by the energy-norm geometry of the solution manifold. The research
task is to formalize the conditioning information flow (the predictor sees the
load token) and to connect the resulting latent geometry to the
identifiability conditions of \cite{klindt2026}.
\end{proposition}

\subsection{Research questions and hypotheses}
\label{sec:rq}

\textbf{RQ1 (anchor quality).} Does a discretization-exact variational anchor
deliver the gains that strong-form FD residuals failed to deliver in
PI-JEPA? \emph{Hypothesis:} yes for elliptic solid statics; the gradient
identity of Lemma~\ref{lem:energy} removes the approximation-mismatch failure
mode, and ablating the anchor will hurt, reversing PI-JEPA's $-8.1\%$
finding. This is the single most falsifiable claim in the proposal.

\textbf{RQ2 (label efficiency).} With pretraining on $10^4$--$10^5$ unlabeled
instances, how many labeled solves are needed to match a fully supervised
Transolver/MeshGraphNets trained on $10^3$--$10^4$ solves? \emph{Target:}
$\ge 10\times$ label reduction at equal relative-$L^2$ error on stress.

\textbf{RQ3 (discretization invariance).} Does
$\mathcal{L}_{\mathrm{inv}}$ produce latents and predictions stable under
mesh refinement/coarsening at test time (train-coarse/test-fine gap), beating
operator baselines on cross-resolution transfer?

\textbf{RQ4 (credit assignment).} How does the benefit decompose across
JEPA prediction, energy anchor, SIGReg, and cross-mesh invariance? (Directly
addresses the ``architecture vs.\ pretraining'' confound PI-JEPA reported.)

\textbf{RQ5 (latent semantics).} Do latents support zero-/few-shot
engineering queries---max von Mises stress, compliance, critical-region
localization---via linear probes, and does latent geometry correlate with
energy-norm distances as Proposition~\ref{prop:collapse} predicts?

\subsection{Novelty deltas relative to closest work}

Relative to PI-JEPA \cite{yee2026pijepa}: unstructured FE meshes and
geometry variation (vs.\ fixed grids and parameter-field variation);
assembled variational anchor with an exactness guarantee (vs.\ FD strong-form
found neutral); elliptic statics and energy functionals (vs.\ parabolic
time-stepping and operator splitting); SIGReg-symmetric training (vs.\ EMA +
ad-hoc regularizer); two theory results. Relative to AeroJEPA: label-free
input-side pretraining (vs.\ pretraining on solved fields); structural FEA;
physics anchoring; mesh-refinement augmentation. Relative to Deep Ritz/DEM:
amortization across an instance distribution with a representation-learning
objective (vs.\ per-instance optimization). Relative to Shape and geometry
SSL: the latents are tied to mechanics, not only to shape statistics.

\section{Execution Plan}
\label{sec:plan}

The plan is organized as a 12-month, four-phase program with explicit
go/no-go gates, designed falsification-first: the cheapest experiments that
could kill the idea run earliest. A single researcher with periodic
collaboration suffices for Phases 0--2; Phase 3 benefits from a second pair
of hands on theory or on the nonlinear solver stack.

\subsection{Phase 0 (weeks 1--6): infrastructure and data}

\textbf{Deliverable: \texttt{fejepa} v0.1 + frozen datasets + reproduced baselines.}
Build the data pipeline: (i) a parametric 2D generator (plates with
holes/notches/inclusions; linear elasticity; randomized Dirichlet/Neumann
programs and load batteries) producing per-instance \texttt{(mesh, $K$, $F$
battery)} archives via Gmsh + FEniCSx or scikit-fem, with sparse $K$ stored
in CSR; (ii) ingestion of external corpora for later phases (SimJEB meshes
and load cases \cite{whalen2021simjeb}; Geo-FNO elasticity
\cite{li2023geofno}; a CAD slice of Thingi10K/Fusion360 for unlabeled-only
pretraining); (iii) a labeled evaluation split solved once with CalculiX/
FEniCSx (these are the \emph{only} solves in the project's training economy;
budget $\sim$5k 2D solves and $\sim$1--2k 3D solves). Implement the encoder
(graph-transformer with physics tokens; Transolver-style slicing as a
variant), $g_\psi$, SIGReg (reference implementation exists), and the
assembly-backed $\Pi_h$ loss as a sparse-matvec autograd op. Reproduce two
supervised baselines (MeshGraphNets-style GNN; Transolver) on the 2D set to
calibrate the evaluation harness. \emph{Gate G0:} energy-anchor unit test ---
training $g_\psi$ alone on single instances must converge to the FE solution
in energy norm (this just re-derives a neural linear solver, but it validates
the entire differentiable-assembly path and Lemma~\ref{lem:energy}
numerically).

\subsection{Phase 1 (months 2--4): core FE-JEPA in 2D + falsification battery}

\textbf{Deliverable: full method on 2D elasticity; ablation table; decision memo.}
Train Eq.~\eqref{eq:loss} on $\ge 3\times 10^4$ unlabeled 2D instances;
fine-tune with $\{16, 64, 256, 1024\}$ labels; evaluate against baselines
trained from scratch on the same label budgets. The falsification battery
(Table~\ref{tab:falsify}) is designed so that each headline claim has a
cheap, pre-registered kill test, mirroring the discipline of running
adversarial experiments before writing.

\begin{table}[h]
\centering
\small
\begin{tabular}{@{}llp{6.6cm}@{}}
\toprule
\textbf{ID} & \textbf{Tests} & \textbf{Kill condition (pre-registered)}\\
\midrule
E1 & RQ1: anchor value & Removing $\mathcal{L}_{\mathrm{phys}}$ changes
fine-tuned error by $<3\%$ relative at every label budget
$\Rightarrow$ anchor is neutral here too; pivot to anchor-only framing
(amortized Ritz) and report the negative JEPA-side result honestly.\\
E2 & RQ4: JEPA value & Anchor-only training (no $\mathcal{L}_{\mathrm{pred}}$)
matches full FE-JEPA within 3\% $\Rightarrow$ SSL component not pulling
weight; investigate masking design before concluding.\\
E3 & collapse & With $\lambda_S{=}0$ and no EMA, latents collapse
(effective rank $\to$ small) and Prop.~\ref{prop:collapse}'s prediction
fails empirically $\Rightarrow$ revisit conditioning paths for load tokens.\\
E4 & RQ3: mesh views & $\mathcal{L}_{\mathrm{inv}}$ off vs.\ on shows no
cross-resolution transfer gap difference $\Rightarrow$ drop the
augmentation claim.\\
E5 & sanity vs.\ naive & Fails to beat SimJEB-style naive polynomial
surrogate \cite{whalen2021simjeb} at any budget $\Rightarrow$ implementation
bug hunt before any further claims.\\
\bottomrule
\end{tabular}
\caption{Pre-registered falsification battery (Phase 1).}
\label{tab:falsify}
\end{table}

\emph{Gate G1 (go/no-go):} proceed to 3D iff (a) E5 passes, (b) at least one
of E1/E2 shows a $\ge 10\%$ relative improvement attributable to its
component at the 64-label budget, and (c) pretraining beats from-scratch at
$\le 256$ labels. Otherwise publish the negative/diagnostic study (still
valuable given PI-JEPA's open question) and stop.

\subsection{Phase 2 (months 4--8): 3D, SimJEB, and external baselines}

\textbf{Deliverable: headline results; preprint draft.}
Scale to 3D tetrahedral meshes; pretrain on unlabeled CAD-derived instances
plus SimJEB geometries (labels held out); fine-tune and evaluate on SimJEB's
official splits (displacement components, magnitude, von Mises stress, four
load cases) and on Geo-FNO elasticity. Baselines: supervised MeshGraphNets
\cite{pfaff2021}, Transolver \cite{wu2024transolver}, Geo-FNO
\cite{li2023geofno}, PINO-style physics-informed fine-tuning
\cite{li2024pino}, a masked-reconstruction (MAE-style) pretraining control
with identical encoder (isolating ``JEPA vs.\ reconstruction''), and a
Shape-style geometry-only SSL control (isolating ``mechanics anchor vs.\
shape statistics''). Report the metric suite of
Section~\ref{sec:metrics} with seeds $\ge 3$ and confidence intervals.
\emph{Gate G2:} headline table shows label-efficiency gains $\ge 5\times$ at
matched error on at least one 3D benchmark, or cross-resolution transfer
wins; either suffices for a strong empirical paper.

\subsection{Phase 3 (months 7--10, overlapping): nonlinear extension and theory}

\textbf{Deliverable: hyperelastic results + formalized Prop.~\ref{prop:collapse} + identifiability connection.}
Replace $\Pi_h$ with a Neo-Hookean stored-energy functional (assembled
elementwise; gradients via autodiff through the element loop or via assembled
residual vectors); evaluate on DeformingPlate-style quasi-static benchmarks.
Honest scope: nonconvexity voids the global-minimizer half of
Lemma~\ref{lem:energy} (stationarity remains), and stabilization
(load-stepping curricula in the loss) is expected to matter---this is
flagged as the project's main technical risk, not hidden. In parallel,
formalize Proposition~\ref{prop:collapse} (Lipschitz decoder assumptions,
conditioning information flow, connection to the isotropy conditions of
\cite{klindt2026}); this theory thread is deliberately matched to the
proposer's strengths and can proceed independently of GPU queues.

\subsection{Phase 4 (months 10--12): writing, release, submission}

Paper, code release (\texttt{fejepa} with frozen configs, data generators,
and the labeled splits' provenance), and a model card stating verified vs.\
unverified claims. Venue ladder, to be confirmed against actual deadlines at
submission time: ICML/NeurIPS/ICLR main track (method + theory framing);
AAAI-27 if Phase-2 results land early enough for its earlier deadline;
CMAME or CMES as the mechanics-community version with extended engineering
validation. The theory results alone (exact anchor; collapse geometry) could
seed a short workshop paper earlier as a flag-plant if scooping risk
materializes.

\subsection{Evaluation protocol and metrics}
\label{sec:metrics}

Primary: relative $L^2$ error on displacement and von Mises stress fields;
\emph{energy gap} $\Pi_h(\hat u)-\Pi_h(U^\star)$ (by Lemma~\ref{lem:energy},
exactly $\tfrac12\norm{\hat u-U^\star}_K^2$; the FE-native metric);
label-efficiency curves (error vs.\ number of labeled solves, with area-under-curve
summary). Secondary: max-von-Mises relative error and critical-element
localization accuracy (the quantities engineers act on); cross-resolution
gap (train-coarse/test-fine); OOD geometry splits (held-out shape families);
wall-clock speedup vs.\ direct solve at matched accuracy; latent diagnostics
(effective rank, SIGReg statistic, probe $R^2$ for compliance). All
comparisons at matched parameter counts and tuned baselines; seeds $\ge 3$;
no metric reported without its definition in the released harness.

\subsection{Compute and cost envelope}

2D phases run on a single modern GPU (24--48\,GB): encoder $\le 30$M params,
sparse anchor adds $O(\mathrm{nnz})$ per step. Phase 2 wants 4--8 GPUs for
$\sim$3--4 weeks of accumulated training (pretraining epochs over $10^5$
instances dominate); total project envelope on the order of 5--8k GPU-hours
--- modest by SSL standards because the anchor replaces, rather than adds to,
expensive labels. Data generation is CPU-bound and parallel (meshing +
assembly; the $\sim$6--7k evaluation solves are a one-time cost).

\subsection{Risk register}

\begin{table}[h]
\centering
\small
\begin{tabular}{@{}p{3.6cm}cp{7.6cm}@{}}
\toprule
\textbf{Risk} & \textbf{Severity} & \textbf{Mitigation}\\
\midrule
Scooped on ``JEPA for FEA'' & High & Speed (Phases 0--1 in $\le$4 months);
theory deltas survive scooping; workshop flag-plant option.\\
Anchor neutral (PI-JEPA repeat) & Medium & Lemma~\ref{lem:energy} makes
neutrality diagnosable (decoder capacity or conditioning, not signal
quality); E1 pre-registered; negative result still answers PI-JEPA's open
question.\\
JEPA adds little over anchor & Medium & E2 isolates it; pivot story
(``amortized Ritz pretraining'') remains novel and useful.\\
Nonlinear instability (Phase 3) & Medium & Load-stepping curricula;
stationarity-based losses; scope paper to linear+hyperelastic only if
needed.\\
3D memory blow-up & Low--Med & Token-slicing encoders (Transolver-class);
region-wise decoding; mixed precision; coarse-mesh pretraining with
fine-mesh invariance views.\\
Benchmark licensing/availability drift & Low & SimJEB is public
(Harvard Dataverse); in-house generator guarantees a fallback.\\
Over-claiming & Low & Pre-registered kill conditions; verified-reference
policy; explicit ``unverified'' flags in the paper for any external claim
not reproduced.\\
\bottomrule
\end{tabular}
\caption{Risk register with mitigations.}
\label{tab:risk}
\end{table}

\subsection{Timeline summary}

\begin{table}[h]
\centering
\small
\begin{tabular}{@{}lp{5.2cm}p{5.6cm}@{}}
\toprule
\textbf{Phase (months)} & \textbf{Output} & \textbf{Gate}\\
\midrule
0 (0--1.5) & pipeline, datasets, baselines, anchor unit test & G0: neural-solver sanity passes\\
1 (2--4) & 2D FE-JEPA + falsification battery & G1: component value $\ge$10\%; beats scratch at $\le$256 labels\\
2 (4--8) & 3D + SimJEB headline table, preprint & G2: $\ge 5\times$ label efficiency or resolution-transfer win\\
3 (7--10) & hyperelastic anchor; formal theory & theory write-up complete\\
4 (10--12) & paper, release, submission & venue deadline (verify dates)\\
\bottomrule
\end{tabular}
\caption{Phased timeline with go/no-go gates.}
\label{tab:timeline}
\end{table}

\section{Expected Contributions}
\label{sec:contrib}

(1) The first JEPA framework for structural FEA on unstructured meshes, with
label-free pretraining over geometry/BC/load space. (2) The assembled-energy
anchor with an exactness lemma showing label-free and supervised energy-norm
training coincide in gradient for linear problems---a clean answer to the
open question raised by PI-JEPA's neutral physics term. (3) Mesh refinement
as a provably exact SSL augmentation, with cross-resolution transfer
evidence. (4) A collapse-geometry result linking energy gaps to latent
separation, connected to current JEPA identifiability theory. (5) An
open-source, falsification-first benchmark harness (generator + frozen
splits + metrics) for label-efficient FEA surrogacy, where the field
currently lacks a standard.

\section{Honest Limitations and Open Issues}
\label{sec:limits}

The exactness guarantee is specific to linear self-adjoint problems;
hyperelasticity retains a variational anchor but loses convexity, and
plasticity/contact require incremental variational formulations that are
explicitly out of the 12-month scope. The anchor supervises toward the FE
solution \emph{on the given mesh}, inheriting its discretization error; the
cross-mesh invariance term mitigates but does not eliminate this, and claims
will be phrased as ``matches FEA,'' not ``matches the continuum.'' Latent
prediction on input-side data may, as in PI-JEPA, contribute less than
architecture or anchor; the experimental design treats this as a result to
be measured, not a risk to be obscured. Several 2026 preprints surveyed here
(AeroJEPA, Shape, LGS) are cited by arXiv identifier with author lists
deliberately left unrecorded pending verification, consistent with a
verified-references-only policy; they must be checked before any of this
material enters a submission. Venue deadlines stated nowhere in this
document should be taken from official CFPs at decision time.

\begin{thebibliography}{99}
\small

\bibitem{raissi2019}
M.~Raissi, P.~Perdikaris, G.~E. Karniadakis.
Physics-informed neural networks: A deep learning framework for solving
forward and inverse problems involving nonlinear partial differential
equations. \emph{Journal of Computational Physics} 378:686--707, 2019.

\bibitem{e2018ritz}
W.~E, B.~Yu.
The Deep Ritz method: A deep learning-based numerical algorithm for solving
variational problems. \emph{Communications in Mathematics and Statistics}
6(1):1--12, 2018.

\bibitem{samaniego2020}
E.~Samaniego, C.~Anitescu, S.~Goswami, V.~M. Nguyen-Thanh, H.~Guo,
K.~Hamdia, X.~Zhuang, T.~Rabczuk.
An energy approach to the solution of partial differential equations in
computational mechanics via machine learning: Concepts, implementation and
applications. \emph{Computer Methods in Applied Mechanics and Engineering}
362:112790, 2020.

\bibitem{nguyen2020dem}
V.~M. Nguyen-Thanh, X.~Zhuang, T.~Rabczuk.
A deep energy method for finite deformation hyperelasticity.
\emph{European Journal of Mechanics --- A/Solids}, 2020.

\bibitem{krishnapriyan2021}
A.~Krishnapriyan, A.~Gholami, S.~Zhe, R.~Kirby, M.~W. Mahoney.
Characterizing possible failure modes in physics-informed neural networks.
\emph{NeurIPS}, 2021.

\bibitem{wang2022ntk}
S.~Wang, X.~Yu, P.~Perdikaris.
When and why PINNs fail to train: A neural tangent kernel perspective.
\emph{Journal of Computational Physics}, 2022. arXiv:2007.14527.

\bibitem{pantidis2023ifenn}
P.~Pantidis, M.~E. Mobasher.
Integrated finite element neural network (I-FENN) for non-local continuum
damage mechanics. \emph{Computer Methods in Applied Mechanics and
Engineering} 404:115766, 2023.

\bibitem{lu2021deeponet}
L.~Lu, P.~Jin, G.~Pang, Z.~Zhang, G.~E. Karniadakis.
Learning nonlinear operators via DeepONet based on the universal
approximation theorem of operators. \emph{Nature Machine Intelligence}
3:218--229, 2021.

\bibitem{li2021fno}
Z.~Li, N.~Kovachki, K.~Azizzadenesheli, B.~Liu, K.~Bhattacharya, A.~Stuart,
A.~Anandkumar.
Fourier neural operator for parametric partial differential equations.
\emph{ICLR}, 2021. arXiv:2010.08895.

\bibitem{li2023geofno}
Z.~Li, D.~Z. Huang, B.~Liu, A.~Anandkumar.
Fourier neural operator with learned deformations for PDEs on general
geometries. \emph{Journal of Machine Learning Research}, 2023.
arXiv:2207.05209.

\bibitem{kovachki2023}
N.~Kovachki, Z.~Li, B.~Liu, K.~Azizzadenesheli, K.~Bhattacharya, A.~Stuart,
A.~Anandkumar.
Neural operator: Learning maps between function spaces with applications to
PDEs. \emph{Journal of Machine Learning Research} 24(89):1--97, 2023.

\bibitem{li2024pino}
Z.~Li, H.~Zheng, N.~Kovachki, D.~Jin, H.~Chen, B.~Liu, K.~Azizzadenesheli,
A.~Anandkumar.
Physics-informed neural operator for learning partial differential
equations. arXiv:2111.03794, 2021.

\bibitem{pfaff2021}
T.~Pfaff, M.~Fortunato, A.~Sanchez-Gonzalez, P.~Battaglia.
Learning mesh-based simulation with graph networks. \emph{ICLR}, 2021.
arXiv:2010.03409.

\bibitem{wu2024transolver}
H.~Wu, H.~Luo, H.~Wang, J.~Wang, M.~Long.
Transolver: A fast transformer solver for PDEs on general geometries.
\emph{ICML}, 2024. arXiv:2402.02366.

\bibitem{whalen2021simjeb}
E.~Whalen, A.~Beyene, C.~Mueller.
SimJEB: Simulated Jet Engine Bracket Dataset.
\emph{Computer Graphics Forum (SGP)}, 2021. arXiv:2105.03534.

\bibitem{takamoto2022pdebench}
M.~Takamoto et al.
PDEBench: An extensive benchmark for scientific machine learning.
\emph{NeurIPS Datasets and Benchmarks}, 2022. arXiv:2210.07182.

\bibitem{mccabe2023mpp}
M.~McCabe et al.
Multiple physics pretraining for physical surrogate models.
arXiv:2310.02994, 2023.

\bibitem{herde2024poseidon}
M.~Herde et al.
Poseidon: Efficient foundation models for PDEs. \emph{NeurIPS}, 2024.
arXiv:2405.19101.

\bibitem{hao2024dpot}
Z.~Hao et al.
DPOT: Auto-regressive denoising operator transformer for large-scale PDE
pre-training. \emph{ICML}, 2024. arXiv:2403.03542.

\bibitem{lecun2022}
Y.~LeCun.
A path towards autonomous machine intelligence (version 0.9.2).
OpenReview position paper, 2022.

\bibitem{assran2023ijepa}
M.~Assran, Q.~Duval, I.~Misra, P.~Bojanowski, P.~Vincent, M.~Rabbat,
Y.~LeCun, N.~Ballas.
Self-supervised learning from images with a joint-embedding predictive
architecture. \emph{CVPR}, 2023. arXiv:2301.08243.

\bibitem{bardes2024vjepa}
A.~Bardes et al.
Revisiting feature prediction for learning visual representations from
video (V-JEPA). 2024. arXiv:2404.08471.

\bibitem{meta2025vjepa2}
M.~Assran et al. (Meta FAIR).
V-JEPA 2: Self-supervised video models enable understanding, prediction and
planning. arXiv:2506.09985, 2025.

\bibitem{bardes2022vicreg}
A.~Bardes, J.~Ponce, Y.~LeCun.
VICReg: Variance-invariance-covariance regularization for self-supervised
learning. \emph{ICLR}, 2022. arXiv:2105.04906.

\bibitem{balestriero2025lejepa}
R.~Balestriero, Y.~LeCun.
LeJEPA: Provable and scalable self-supervised learning without the
heuristics. arXiv:2511.08544, 2025.

\bibitem{klindt2026}
D.~Klindt, Y.~LeCun, R.~Balestriero.
When does LeJEPA learn a world model? arXiv:2605.26379, 2026.

\bibitem{yee2026pijepa}
B.~Yee, P.~Koh.
PI-JEPA: Closing the data asymmetry in subsurface surrogate modeling via
label-free pretraining. arXiv:2604.01349, 2026.

\end{thebibliography}

\end{document}
